\section{Extended Cloud and Edge Literature}
\label{app:extended-literature}
\subsection{Cloud, Microservice, and Load-Balancing Literature}

Recent cloud resource-management studies emphasize that adaptive schedulers
must operate under workload uncertainty, heterogeneous resources, and competing
objectives such as latency, throughput, energy, and utilization
\cite{Khan2022MLCloudReview}. Reviews of DRL-based cloud control identify the
same practical issues that matter here: reward design, non-stationary demand,
and operational deployment \cite{Zhou2024DRLCloudReview}. DRL has been used
for algorithm selection in hierarchical cloud scheduling
\cite{Zhou2022DRLSelector}, feedback-control-based cloud resource allocation
\cite{Chen2022RLFeedback}, multi-objective virtual-machine placement
\cite{Caviglione2021VMPlacement}, and multidimensional allocation in
distributed data centers \cite{Wei2023TNSM}. These works show how learned
control can account for delayed effects and competing objectives.
However, their controlled objects are usually tasks, virtual machines, jobs, or
data-center resources. Our controlled objects are lower in the messaging
stack: subscriptions, persistent client connections, and observed subscriber
pods.

Cloud scheduling offers several distinct approaches, not all of them
learning-based. Belgacem et al. use a spacing multi-objective antlion algorithm
to optimize makespan and virtual-machine cost \cite{Belgacem2020DynamicResource}.
Ding et al. instead use Q-learning for energy-efficient task scheduling
\cite{Ding2020QLearningCloud}, while Chen et al. use actor--critic DRL for
resource allocation in cloud data centers \cite{Chen2022ActorCritic}.
Collaborative scheduling addresses heterogeneous workflows
\cite{Chen2023CollaborativeScheduling}; workload-time windows provide temporal
context for cloud-service allocation \cite{Chen2023WorkloadWindows}.
Li et al. apply DRL to energy-aware Spark scheduling \cite{Li2023SparkDRL}.
Sanjalawe et al. provide a comprehensive review of AI-driven cloud job
scheduling \cite{Sanjalawe2025AIJobScheduling}.

The controlled workload also matters. Dong et al. use DQN for manufacturing
task scheduling \cite{Dong2020CloudManufacturing}, and Chen et al. study
attention-based scheduling across manufacturing cloud and edge resources
\cite{Chen2023CloudEdgeManufacturing}. Cheng et al. introduce preemption into
cost-aware cloud scheduling \cite{Cheng2023PreemptiveCloud}; Liu et al. consider
inter-region bandwidth and traffic costs for data-intensive jobs
\cite{Liu2026GeoDRL}. These papers supply scheduling context. Our problem
instead requires subscription reassociation while retaining queue identity
and delivery ownership across persistent connections.

Edge research is relevant here in two narrower respects: decentralized
observations and explicit tradeoffs among resource costs. Iftikhar et al.
provide a taxonomy of AI-based fog and edge computing \cite{Iftikhar2023EdgeAI},
while Hortelano et al. survey RL-based computation offloading
\cite{Hortelano2023RLOffloading}. A concrete decentralized example is Nieto
et al.'s per-device actor--critic offloading policy, which trades battery-related
QoE against task-latency requirements \cite{Nieto2024DRLOffloading}. Hou et al.
explicitly formulate distributed offloading with partial global information
\cite{Hou2023PODRL}. Partial observability is therefore already addressed in
earlier work. It is not, by itself, a new contribution of our controller.

Other offloading papers inform the choice of cost terms, rather than the
transport handoff mechanism. Xu et al. optimize cloud-edge resource rental and
allocation costs \cite{Xu2022CloudEdgeResource}. Huang et al. combine
computation offloading with resource allocation in wireless-powered edge
networks \cite{Huang2020Offloading}. Pan et al. model dependent tasks and train
an offloading scheduler using a modified PPO method \cite{Pan2023ICNOffloading}.
Jayanetti et al. consider energy and execution time for precedence-constrained
tasks \cite{Jayanetti2022EdgeCloudScheduling}. Mitsis et al. incorporate price
and risk into offloading decisions \cite{Mitsis2022RiskOffloading}, whereas Cui
et al. optimize caching and request routing for latency-sensitive content
distribution \cite{Cui2023IoTEdgeCloud}. Our migration and connection costs
likewise make operational tradeoffs explicit, but their definitions and the
subscription-ownership contract are specific to persistent messaging.

Microservice autoscaling provides a closer infrastructure comparison.
Abdel Khaleq and Ra use reinforcement learning to control horizontal scaling
from resource demand \cite{AbdelKhaleq2021Autoscaling}. HephaestusForge learns
microservice placement across Kubernetes clusters with multiple deployment
objectives \cite{Santos2025HephaestusForge}. These approaches change service
capacity or placement. Our mechanism addresses the complementary client-side
step: existing subscription-bearing sockets must reach the added capacity,
and candidate promotion must account for both subscription and backlog pressure.

